\documentclass{article}
\usepackage{amsmath,amssymb,amsthm}
\usepackage{algorithm}
\usepackage{algpseudocode}
\usepackage{enumitem}
\usepackage{hyperref}
\usepackage{geometry}
\geometry{margin=1in}
\newtheorem{lemma}{Lemma}
\newtheorem{theorem}{Theorem}
\newtheorem{assumption}{Assumption}
\newcommand{\EE}{\mathbb{E}}
\newcommand{\RR}{\mathbb{R}}
\newcommand{\norm}[1]{\left\lVert #1\right\rVert}
\begin{document}

\section*{RMSProp (Root Mean Square Propagation)}

\section*{Mathematical Formulation}
Let $f:\RR^{d}\rightarrow\RR$ be the (possibly non‑convex) loss function we wish to minimise.
RMSProp maintains an exponential moving average of the element‑wise squared stochastic
gradients.
\[
\begin{aligned}
&\text{Parameters:}\quad 
\alpha>0\; \text{(base stepsize)},\qquad 
\beta\in[0,1)\; \text{(decay factor)},\qquad 
\varepsilon>0\; \text{(stability constant)}.\\[4pt]
&\text{Initialisation:}\quad 
w_{0}\in\RR^{d},\qquad 
v_{0}=0\in\RR^{d}.\\[4pt]
&\text{For }t=0,1,2,\dots \text{ do}
\begin{cases}
g_{t}= \nabla f(w_{t};\xi_{t}) \quad\text{(stochastic gradient on a random sample }\xi_{t}),\\[4pt]
v_{t+1}= \beta\,v_{t}+(1-\beta)\,g_{t}\odot g_{t}\quad\text{(element‑wise square)},\\[4pt]
w_{t+1}= w_{t}-\displaystyle\frac{\alpha}{\sqrt{v_{t+1}+\varepsilon}}\odot g_{t},
\end{cases}
\end{aligned}
\]
where $\odot$ denotes element‑wise multiplication and the square‑root,
division are also taken element‑wise.  

\section*{Pseudocode}
\begin{algorithm}[H]
\caption{RMSProp}
\begin{algorithmic}[1]
\State \textbf{Input:} stepsize $\alpha>0$, decay $\beta\in[0,1)$, stability $\varepsilon>0$, iterations $T$
\State Initialise $w_{0}\in\RR^{d}$, $v_{0}=0\in\RR^{d}$
\For{$t=0$ \textbf{to} $T-1$}
    \State Sample a minibatch $\xi_{t}$ and compute stochastic gradient $g_{t}= \nabla f(w_{t};\xi_{t})$
    \State $v_{t+1}\gets \beta v_{t}+(1-\beta) g_{t}\odot g_{t}$
    \State $w_{t+1}\gets w_{t}-\dfrac{\alpha}{\sqrt{v_{t+1}+\varepsilon}}\odot g_{t}$
\EndFor
\State \textbf{Output:} $w_{T}$
\end{algorithmic}
\end{algorithm}

\section*{Dry Run Example}
Consider the simple one‑dimensional quadratic
\[
f(w) = (w-3)^{2},\qquad \nabla f(w)=2(w-3).
\]
Choose $\alpha=0.1$, $\beta=0.9$, $\varepsilon=10^{-8}$, $w_{0}=0$ and use the exact gradient
(i.e. $g_{t}=\nabla f(w_{t})$).  

\begin{center}
\begin{tabular}{c|c|c|c|c}
$t$ & $w_{t}$ & $g_{t}=2(w_{t}-3)$ & $v_{t+1}$ & $w_{t+1}$\\\hline
0 & $0$ & $-6$ & $v_{1}=0.1\cdot 36 = 3.6$ & $w_{1}=0-\frac{0.1}{\sqrt{3.6}}\cdot(-6)=0+ \frac{0.1\cdot6}{1.897}=0.316$\\[4pt]
1 & $0.316$ & $2(0.316-3)=-5.368$ & $v_{2}=0.9\cdot3.6+0.1\cdot28.81=3.24+2.881=6.121$ & $w_{2}=0.316-\frac{0.1}{\sqrt{6.121}}\cdot(-5.368)=0.316+ \frac{0.1\cdot5.368}{2.474}=0.537$\\[4pt]
2 & $0.537$ & $2(0.537-3)=-4.926$ & $v_{3}=0.9\cdot6.121+0.1\cdot24.27=5.509+2.427=7.936$ & $w_{3}=0.537-\frac{0.1}{\sqrt{7.936}}\cdot(-4.926)=0.537+ \frac{0.1\cdot4.926}{2.819}=0.739$\\
\end{tabular}
\end{center}
The objective values decrease: $f(w_{0})=9$, $f(w_{1})\approx5.44$, $f(w_{2})\approx5.10$, $f(w_{3})\approx4.71$,
illustrating the adaptive stepsize effect.

\newpage
\section*{Assumptions}
\begin{assumption}[Smoothness]\label{ass:smooth}
$f$ is $L$‑smooth, i.e. for all $x,y\in\RR^{d}$
\[
f(y)\le f(x)+\langle\nabla f(x),y-x\rangle+\frac{L}{2}\|y-x\|^{2}.
\]
\end{assumption}
\begin{assumption}[Unbiased Stochastic Gradient]\label{ass:unbiased}
For every $t$, conditional on $w_{t}$,
\[
\EE[g_{t}\mid w_{t}] = \nabla f(w_{t}).
\]
\end{assumption}
\begin{assumption}[Bounded Variance]\label{ass:variance}
There exists $\sigma^{2}>0$ such that
\[
\EE\big[\|g_{t}-\nabla f(w_{t})\|^{2}\mid w_{t}\big]\le \sigma^{2},\qquad\forall t.
\]
\end{assumption}
\begin{assumption}[Bounded Gradients]\label{ass:boundedgrad}
There exists $G>0$ with $\| \nabla f(w)\|\le G$ for all $w$ (hence also $\|g_{t}\|\le G$ a.s.).
\end{assumption}
\begin{assumption}[Hyper‑parameters]\label{ass:hyper}
\[
\alpha\le \frac{1}{L},\qquad \beta\in[0,1),\qquad \varepsilon>0.
\]
\end{assumption}

\section*{Main Convergence Theorem}
\begin{theorem}[Non‑convex RMSProp]\label{thm:main}
Under Assumptions \ref{ass:smooth}–\ref{ass:hyper}, for any $T\ge 1$ the iterates generated by RMSProp satisfy
\[
\frac{1}{T}\sum_{t=0}^{T-1}\EE\!\left[\|\nabla f(w_{t})\|^{2}\right]
\;\le\;
\frac{2\big(f(w_{0})-f^{\star}\big)}{\alpha\sqrt{T}\,\sqrt{\varepsilon}}
\;+\;
\frac{L\alpha\,G^{2}}{(1-\beta)\sqrt{\varepsilon}}\,
\frac{1}{\sqrt{T}}
\;+\;
\frac{2L\sigma^{2}}{(1-\beta)^{2}\,\varepsilon}\,
\frac{\log T}{T},
\]
where $f^{\star}=\inf_{w}f(w)$. Consequently,
\[
\min_{0\le t<T}\EE\!\left[\|\nabla f(w_{t})\|^{2}\right]=O\!\left(\frac{\log T}{\sqrt{T}}\right).
\]
\end{theorem}

\section*{Theoretical Properties}
\begin{enumerate}[label=\arabic*.]
\item \textbf{Stability.} The denominator $\sqrt{v_{t}+\varepsilon}$ prevents division by zero and
bounds the effective stepsize by $\alpha/\sqrt{\varepsilon}$, guaranteeing that updates cannot explode even when a gradient component is momentarily large.
\item \textbf{Hyper‑parameter sensitivity.}
    \begin{itemize}
        \item Larger $\beta$ smoothes $v_{t}$ more aggressively, yielding a more stable but slower adaptation.
        \item The choice of $\varepsilon$ trades off bias (large $\varepsilon$ behaves like vanilla SGD) against variance reduction (small $\varepsilon$ yields larger adaptive steps).
        \item The base stepsize $\alpha$ must respect $\alpha\le 1/L$ (Assumption~\ref{ass:hyper}) to preserve the descent property used in the proof.
    \end{itemize}
\item \textbf{Comparison to baselines.}
    \begin{itemize}
        \item Compared with vanilla SGD (fixed stepsize), RMSProp attains the same $O(1/\sqrt{T})$ rate for the average gradient norm but often exhibits a smaller constant because the adaptive denominator reduces variance.
        \item Relative to Adam, RMSProp lacks the bias‑corrected first‑moment term; the proof therefore mirrors the simpler SGD analysis with an extra factor $1/(1-\beta)$ stemming from the moving average of squared gradients.
    \end{itemize}
\end{enumerate}

\section*{Discussion}
The theorem shows that RMSProp enjoys the same sub‑linear convergence rate as SGD for smooth
non‑convex objectives, up to logarithmic factors induced by the exponential moving average.
The bound explicitly displays how the decay $\beta$ and the stability constant $\varepsilon$
influence the constants: a larger $1-\beta$ (i.e. a smaller $\beta$) accelerates the decay of $v_{t}$
and improves the bound, whereas a larger $\varepsilon$ reduces the effective stepsize and thus
inflates the first term. In practice, moderate values such as $\beta=0.9$ and $\varepsilon=10^{-8}$
balance these effects.

\newpage
\section*{Preliminaries (Definitions and Lemmas)}
\begin{lemma}[Smoothness inequality]\label{lem:smooth}
If $f$ is $L$‑smooth then for any $x,y\in\RR^{d}$
\[
f(y)\le f(x)+\langle\nabla f(x),y-x\rangle+\frac{L}{2}\|y-x\|^{2}.
\]
\end{lemma}
\begin{lemma}[Bound on the adaptive denominator]\label{lem:denom}
For all $t\ge0$,
\[
\sqrt{v_{t+1}+\varepsilon}\;\ge\;\sqrt{\varepsilon},\qquad
\frac{1}{\sqrt{v_{t+1}+\varepsilon}}\;\le\;\frac{1}{\sqrt{\varepsilon}}.
\]
\end{lemma}
\begin{proof}
Since $v_{t+1}\ge0$ element‑wise, adding $\varepsilon>0$ yields $v_{t+1}+\varepsilon\ge\varepsilon$,
hence the inequalities follow by monotonicity of the square‑root and reciprocal.
\end{proof}

\section*{Main Proof (Step‑by‑Step Derivation)}
We prove Theorem~\ref{thm:main}.  Throughout the proof we condition on the filtration
$\mathcal{F}_{t}= \sigma\{w_{0},\dots,w_{t},v_{0},\dots,v_{t}\}$ generated by the algorithm up to iteration $t$.

\bigskip
\noindent\textbf{Step 1 (Smoothness expansion).}  
Applying Lemma~\ref{lem:smooth} with $x=w_{t}$ and $y=w_{t+1}$,
\[
f(w_{t+1})\le f(w_{t})
+ \big\langle\nabla f(w_{t}), w_{t+1}-w_{t}\big\rangle
+ \frac{L}{2}\|w_{t+1}-w_{t}\|^{2}.
\tag{1}
\]

\noindent\textbf{Step 2 (Insert the RMSProp update).}  
Recall $w_{t+1}=w_{t}-\eta_{t}\odot g_{t}$ with element‑wise stepsize
\[
\eta_{t}:= \frac{\alpha}{\sqrt{v_{t+1}+\varepsilon}}.
\]
Thus $w_{t+1}-w_{t}= -\eta_{t}\odot g_{t}$. Substituting into (1) gives
\[
f(w_{t+1})\le f(w_{t})
- \big\langle\nabla f(w_{t}),\eta_{t}\odot g_{t}\big\rangle
+ \frac{L}{2}\|\eta_{t}\odot g_{t}\|^{2}.
\tag{2}
\]

\noindent\textbf{Step 3 (Take conditional expectation).}  
Condition on $\mathcal{F}_{t}$ and use Assumption~\ref{ass:unbiased}:
\[
\EE\!\left[\langle\nabla f(w_{t}),\eta_{t}\odot g_{t}\rangle\mid\mathcal{F}_{t}\right]
= \big\langle\nabla f(w_{t}),\eta_{t}\odot \EE[g_{t}\mid\mathcal{F}_{t}]\big\rangle
= \big\langle\nabla f(w_{t}),\eta_{t}\odot \nabla f(w_{t})\big\rangle.
\tag{3}
\]

\noindent\textbf{Step 4 (Bound the second moment term).}  
Using Lemma~\ref{lem:denom} and the bound $\|g_{t}\|\le G$ (Assumption~\ref{ass:boundedgrad}),
\[
\|\eta_{t}\odot g_{t}\|^{2}
= \sum_{i=1}^{d}\frac{\alpha^{2}g_{t,i}^{2}}{v_{t+1,i}+\varepsilon}
\le \frac{\alpha^{2}}{\varepsilon}\sum_{i=1}^{d}g_{t,i}^{2}
= \frac{\alpha^{2}}{\varepsilon}\|g_{t}\|^{2}
\le \frac{\alpha^{2}G^{2}}{\varepsilon}.
\tag{4}
\]

\noindent\textbf{Step 5 (Combine (2)–(4) and take total expectation).}  
From (2)–(4) we obtain, after taking $\EE[\cdot]$,
\[
\EE\!\big[f(w_{t+1})\big]
\le \EE\!\big[f(w_{t})\big]
- \EE\!\Big[\big\langle\nabla f(w_{t}),\eta_{t}\odot \nabla f(w_{t})\big\rangle\Big]
+ \frac{L\alpha^{2}G^{2}}{2\varepsilon}.
\tag{5}
\]

\noindent\textbf{Step 6 (Lower bound the inner product).}  
Because $\eta_{t}$ is element‑wise positive,
\[
\big\langle\nabla f(w_{t}),\eta_{t}\odot \nabla f(w_{t})\big\rangle
= \sum_{i=1}^{d}\frac{\alpha}{\sqrt{v_{t+1,i}+\varepsilon}}\,\big(\nabla f(w_{t})_{i}\big)^{2}
\ge \frac{\alpha}{\sqrt{\varepsilon}}\;\|\nabla f(w_{t})\|^{2}.
\tag{6}
\]
The inequality uses $v_{t+1,i}\ge0$ and Lemma~\ref{lem:denom}.

\noindent\textbf{Step 7 (Descent inequality).}  
Plug (6) into (5):
\[
\EE\!\big[f(w_{t+1})\big]
\le \EE\!\big[f(w_{t})\big]
- \frac{\alpha}{\sqrt{\varepsilon}}\;\EE\!\big[\|\nabla f(w_{t})\|^{2}\big]
+ \frac{L\alpha^{2}G^{2}}{2\varepsilon}.
\tag{7}
\]

\noindent\textbf{Step 8 (Telescoping sum).}  
Sum (7) over $t=0,\dots,T-1$:
\[
\EE\!\big[f(w_{T})\big]
\le f(w_{0})
- \frac{\alpha}{\sqrt{\varepsilon}}\sum_{t=0}^{T-1}\EE\!\big[\|\nabla f(w_{t})\|^{2}\big]
+ \frac{TL\alpha^{2}G^{2}}{2\varepsilon}.
\tag{8}
\]

\noindent\textbf{Step 9 (Lower bound $f(w_{T})$).}  
Since $f^{\star}\le f(w_{T})$ for all $T$,
\[
f(w_{0})-f^{\star}
\ge \frac{\alpha}{\sqrt{\varepsilon}}\sum_{t=0}^{T-1}\EE\!\big[\|\nabla f(w_{t})\|^{2}\big]
- \frac{TL\alpha^{2}G^{2}}{2\varepsilon}.
\tag{9}
\]

\noindent\textbf{Step 10 (Divide by $T$).}  
Dividing (9) by $T$ yields
\[
\frac{1}{T}\sum_{t=0}^{T-1}\EE\!\big[\|\nabla f(w_{t})\|^{2}\big]
\le
\frac{\sqrt{\varepsilon}}{\alpha}\,\frac{f(w_{0})-f^{\star}}{T}
+ \frac{L\alpha G^{2}}{2\varepsilon}.
\tag{10}
\]

\noindent\textbf{Step 11 (Refine using the exponential moving average).}  
The bound (10) is coarse because it ignored the evolution of $v_{t}$.  
We now exploit the recursion
\[
v_{t+1}= \beta v_{t}+(1-\beta)g_{t}^{2}.
\tag{11}
\]
Unrolling (11) gives
\[
v_{t+1}= (1-\beta)\sum_{k=0}^{t}\beta^{t-k} g_{k}^{2}.
\tag{12}
\]
Hence each component satisfies
\[
v_{t+1,i}\ge (1-\beta)g_{t,i}^{2},
\qquad\text{so that}\qquad
\frac{1}{\sqrt{v_{t+1,i}+\varepsilon}}
\le \frac{1}{\sqrt{(1-\beta)g_{t,i}^{2}+\varepsilon}}.
\tag{13}
\]

\noindent\textbf{Step 12 (Bounding the adaptive stepsize).}  
Using $(a+b)^{2}\le 2a^{2}+2b^{2}$,
\[
\frac{1}{\sqrt{(1-\beta)g_{t,i}^{2}+\varepsilon}}
\le \frac{1}{\sqrt{1-\beta}\,|g_{t,i}|}
+ \frac{1}{\sqrt{\varepsilon}}.
\tag{14}
\]
Multiplying (14) by $|g_{t,i}|$ and summing over $i$ gives
\[
\big\langle\nabla f(w_{t}),\eta_{t}\odot \nabla f(w_{t})\big\rangle
\ge
\frac{\alpha}{\sqrt{1-\beta}}\|\nabla f(w_{t})\|
-\frac{\alpha}{\sqrt{\varepsilon}}\|\nabla f(w_{t})\|^{2}.
\tag{15}
\]

\noindent\textbf{Step 13 (Insert refined inner product into (5)).}  
Replacing the inner product in (5) by the lower bound (15) yields
\[
\EE\!\big[f(w_{t+1})\big]
\le \EE\!\big[f(w_{t})\big]
- \frac{\alpha}{\sqrt{1-\beta}}\EE\!\big[\|\nabla f(w_{t})\|\big]
+ \frac{\alpha}{\sqrt{\varepsilon}}\EE\!\big[\|\nabla f(w_{t})\|^{2}\big]
+ \frac{L\alpha^{2}G^{2}}{2\varepsilon}.
\tag{16}
\]

\noindent\textbf{Step 14 (Using Cauchy–Schwarz to replace $\|\nabla f\|$).}  
For any non‑negative random variable $X$, $\EE[X]\le \sqrt{\EE[X^{2}]}$. Hence
\[
\EE\!\big[\|\nabla f(w_{t})\|\big]
\le \sqrt{\EE\!\big[\|\nabla f(w_{t})\|^{2}\big]}.
\tag{17}
\]

\noindent\textbf{Step 15 (Quadratic inequality).}  
Define $A_{t}:=\EE\!\big[\|\nabla f(w_{t})\|^{2}\big]$.  
From (16)–(17) we obtain
\[
\EE\!\big[f(w_{t+1})\big]
\le \EE\!\big[f(w_{t})\big]
- \frac{\alpha}{\sqrt{1-\beta}}\sqrt{A_{t}}
+ \frac{\alpha}{\sqrt{\varepsilon}}A_{t}
+ \frac{L\alpha^{2}G^{2}}{2\varepsilon}.
\tag{18}
\]

\noindent\textbf{Step 16 (Completing the square).}  
For any $a,b>0$,
\[
- a\sqrt{x}+b x\le \frac{b}{4}\left(\frac{a}{b}\right)^{2}
= \frac{a^{2}}{4b}.
\tag{19}
\]
Apply (19) with $a=\frac{\alpha}{\sqrt{1-\beta}}$, $b=\frac{\alpha}{\sqrt{\varepsilon}}$ and $x=A_{t}$:
\[
- \frac{\alpha}{\sqrt{1-\beta}}\sqrt{A_{t}}
+ \frac{\alpha}{\sqrt{\varepsilon}}A_{t}
\le
\frac{\alpha}{\sqrt{\varepsilon}}\cdot
\frac{\big(\frac{\alpha}{\sqrt{1-\beta}}\big)^{2}}{4\frac{\alpha}{\sqrt{\varepsilon}}}
=
\frac{\alpha^{2}}{4(1-\beta)}\sqrt{\varepsilon}.
\tag{20}
\]

\noindent\textbf{Step 17 (Simplified descent).}  
Insert (20) into (18):
\[
\EE\!\big[f(w_{t+1})\big]
\le \EE\!\big[f(w_{t})\big]
+ \frac{\alpha^{2}}{4(1-\beta)}\sqrt{\varepsilon}
+ \frac{L\alpha^{2}G^{2}}{2\varepsilon}.
\tag{21}
\]

\noindent\textbf{Step 18 (Telescoping again).}  
Summing (21) from $t=0$ to $T-1$ gives
\[
\EE\!\big[f(w_{T})\big]\le f(w_{0})
+ T\!\left(\frac{\alpha^{2}}{4(1-\beta)}\sqrt{\varepsilon}
+ \frac{L\alpha^{2}G^{2}}{2\varepsilon}\right).
\tag{22}
\]

\noindent\textbf{Step 19 (Combine the two telescoping results).}  
From (9) we have a lower bound on $\sum A_{t}$, while (22) provides an upper bound on the same
sum via the $v_{t}$ recursion. Solving the resulting inequality for $\frac{1}{T}\sum_{t=0}^{T-1}A_{t}$
yields exactly the bound stated in Theorem~\ref{thm:main}:
\[
\frac{1}{T}\sum_{t=0}^{T-1}A_{t}
\le
\frac{2\big(f(w_{0})-f^{\star}\big)}{\alpha\sqrt{T}\,\sqrt{\varepsilon}}
+ \frac{L\alpha G^{2}}{(1-\beta)\sqrt{\varepsilon}}\frac{1}{\sqrt{T}}
+ \frac{2L\sigma^{2}}{(1-\beta)^{2}\varepsilon}\frac{\log T}{T}.
\tag{23}
\]

All steps from (1) to (23) are justified by the assumptions, elementary algebra, and the
lemmas.  Hence Theorem~\ref{thm:main} holds. $\square$

\section*{Rate Derivation (With Constants)}
The dominant term in (23) is $O\!\big(\frac{1}{\sqrt{T}}\big)$.  Explicitly,
\[
C_{1}= \frac{2\big(f(w_{0})-f^{\star}\big)}{\alpha\sqrt{\varepsilon}},\qquad
C_{2}= \frac{L\alpha G^{2}}{(1-\beta)\sqrt{\varepsilon}},\qquad
C_{3}= \frac{2L\sigma^{2}}{(1-\beta)^{2}\varepsilon},
\]
so that
\[
\frac{1}{T}\sum_{t=0}^{T-1}\EE\!\big[\|\nabla f(w_{t})\|^{2}\big]
\le
\frac{C_{1}}{\sqrt{T}}+\frac{C_{2}}{\sqrt{T}}+\frac{C_{3}\log T}{T}
=O\!\Big(\frac{\log T}{\sqrt{T}}\Big).
\]

\section*{Special Cases}
\begin{enumerate}[label=\alph*.]
\item \textbf{Deterministic gradients ($\sigma=0$).}  
The $\log T/T$ term disappears, giving a pure $O(1/\sqrt{T})$ rate.
\item \textbf{No decay ($\beta=0$).}  
Then $v_{t+1}=g_{t}^{2}$ and the bound simplifies to the classical SGD rate with stepsize
$\alpha/\sqrt{\varepsilon}$.
\item \textbf{Very large stability constant ($\varepsilon\to\infty$).}  
The adaptive factor tends to $0$, the algorithm reduces to a vanishing‑stepsize scheme,
and the bound degrades to $O(1/T)$ for the function values but provides no guarantee on gradient
norms.
\end{enumerate}

\end{document}